\begin{explanationbox}
	\n\textbf{\textcolor{CExplanation}{一句话直觉}}\n\n过反比例函数图象上任意一点向两坐标轴作垂线，所得矩形面积恒等于 $|k|$，与点的位置无关。\n\n

	\medskip
	\n\textbf{\textcolor{CExplanation}{核心拆解}}\n
	\begin{description}[style=nextline, leftmargin=2.8em, labelsep=0.5em, itemsep=0.45em]
		\n
		\item[\textbf{要点一（坐标乘积恒定）：}]
			因为点 $P(x_0, y_0)$ 在函数 $y = \dfrac{k}{x}$ 的图象上，所以 $x_0 y_0 = k$。因此 $|x_0 y_0| = |k|$ 恒成立。\n
		\item[\textbf{要点二（面积与坐标关系）：}]
			矩形 $OAPB$ 的边长分别为 $|x_0|$ 和 $|y_0|$，面积 $S = |x_0| \cdot |y_0| = |x_0 y_0| = |k|$。\n
	\end{description}
	\n\n

	\medskip
	\n\textbf{\textcolor{CExplanation}{几何本质（面积定值）}}\n\n反比例函数图象是双曲线，$k$ 的绝对值控制了双曲线与坐标轴之间“开口”宽度。过双曲线上任一点作坐标轴的垂线，所得矩形面积恒为定值 $|k|$。\n\n

	\medskip
	\n\textbf{\textcolor{CExplanation}{代数意义（与点无关）}}\n\n$k$ 一旦确定，矩形面积 $|k|$ 就确定，与所取点的具体位置无关。这一性质可用于直接求 $k$。\n\n

	\medskip
	\n\textbf{\textcolor{CExplanation}{考点价值（题型基础）}}\n
	\begin{itemize}[leftmargin=*, itemsep=0.5em, topsep=0.4em]
		\n
		\item \textbf{考法一（已知面积求 $k$）：} 给出矩形或三角形面积，通过 $S=|k|$ 或 $S_{\triangle}=\frac{1}{2}|k|$ 计算 $|k|$，再结合图象象限确定 $k$ 符号。\n
		\item \textbf{考法二（已知 $k$ 求面积）：} 直接代入公式得面积大小。\n
	\end{itemize}
	\n\n

	\medskip
	\n\textbf{\textcolor{CExplanation}{顿悟点}}\n\n面积与点位置无关，只与 $k$ 的绝对值有关。因此，只要知道一个点的坐标，就能算面积；反过来，知道面积就能求出 $|k|$。\n\n

	\medskip
	\n\textbf{\textcolor{CExplanation}{使用场景}}\n
	\begin{itemize}[leftmargin=*, itemsep=0.5em, topsep=0.4em]
		\n
		\item \textbf{场景一（求函数解析式）：} 已知图象上一点与坐标轴围成的面积，求 $k$ 值。\n
		\item \textbf{场景二（比较性质）：} 通过面积大小比较不同反比例函数 $|k|$ 的大小。\n
	\end{itemize}
	\n
\end{explanationbox}
